\chapter{Normals and Surface Reconstruction}\label{ch:recon}

After registration, the scans form a single point cloud that samples the surface of the object.
Surface reconstruction turns it into a watertight mesh. Almost every method needs oriented
normals, so the chapter starts with normal estimation, orientation and outlier removal, and then
studies the extraction of iso-surfaces and three implicit reconstruction methods: Hoppe's signed
distance, radial basis functions and Poisson reconstruction. It follows Tema 7 and the companion
notebook, which implements every method on a genus-one test surface; the survey of
\citet{Berger2017} covers the field broadly.

\section{Normal estimation}

\subsection{Total least squares and PCA}

A point cloud carries no normals, but the neighbours of a point lie approximately on its tangent
plane. The plane $\{\bx:(\bx-\bc)\cdot\bn=0\}$ that best fits points $\bp_1,\dots,\bp_k$ in the
\emph{total least squares} sense minimises the sum of squared orthogonal distances:
\begin{equation}
\min_{\bc,\ \norm{\bn}=1}\ \sum_{i=1}^k\bigl((\bp_i-\bc)\cdot\bn\bigr)^2 .
\end{equation}
For fixed $\bn$ the optimal $\bc$ is the centroid $\bar\bp$. Then the objective is
$\bn\T\mC\bn$ with the covariance matrix
\begin{equation}
\mC=\sum_{i=1}^k(\bp_i-\bar\bp)(\bp_i-\bar\bp)\T ,
\end{equation}
and minimising a Rayleigh quotient over unit vectors gives the eigenvector of the smallest
eigenvalue (\cref{app:linalg}). Writing the eigenvalues as $\lambda_0\le\lambda_1\le\lambda_2$, the
normal is the eigenvector of $\lambda_0$, the other two span the tangent plane, and $\lambda_0$ is
the residual of the fit. Applied to the $k$ nearest neighbours of every point, this is
\emph{PCA normal estimation}, and it is fully vectorisable: build a kd-tree, gather the neighbours,
compute all $3\times3$ covariances and their eigen-decompositions in batch.

\subsection{The size of the neighbourhood}

The choice of $k$ (or of a radius) is a bias--variance trade-off \citep{Mitra2003}. With noise of
standard deviation $\sigma$, a larger neighbourhood averages more points and the angular error
decreases (roughly like $\sigma/r$ for a neighbourhood of radius $r$). With curvature $\kappa$, a
larger neighbourhood includes points far from the tangent plane and the error grows (roughly like
$\kappa r$). The optimal $k^\star$ balances both and increases with the noise
(\cref{fig:normals-k}).

\fig[0.85\textwidth]{ch09_normals_k.pdf}{PCA normal estimation on a noisy genus-one surface.
Median angle error against the number of neighbours $k$ for several noise levels. Without noise the
error grows with $k$ because of curvature (slope $k^{1/2}$); with noise it first decreases (slope
about $k^{-1}$) and then joins the curvature branch, so the best $k$ grows with
$\sigma$.}{fig:normals-k}

\subsection{Consistent orientation}

PCA returns the normal line, not its direction: the sign of each eigenvector is arbitrary. When the
scanner position is known, every normal is simply flipped to point towards the sensor. Otherwise,
\citet{Hoppe1992} propagate the orientation over a graph:
\begin{enumerate}
  \item build the \emph{Riemannian graph}, connecting each point to its $k$ nearest neighbours;
  \item weight each edge $(i,j)$ by $1-\abs{\bn_i\cdot\bn_j}$, which is small between nearly
  parallel normals;
  \item compute a minimum spanning tree, start from a point whose orientation is known (for
  example the highest point, whose normal must point up) and traverse the tree, flipping $\bn_j$
  whenever $\bn_i\cdot\bn_j<0$ for its parent $i$.
\end{enumerate}
The minimum spanning tree favours propagation across flat regions, where the decision is reliable,
and delays it across sharp edges and thin parts, where it is not. On the test surface the fraction
of outward normals goes from $50.3\%$ (random signs) to $100\%$ (\cref{fig:orientation}).

\fig[0.7\textwidth]{ch09_orientation.png}{Normal orientation by minimum spanning tree
propagation. Left: PCA normals with random signs (red: pointing inwards). Right: after propagation
all normals point outwards.}{fig:orientation}

\subsection{Outliers and what the eigenvalues say}

Isolated points far from the surface corrupt normals and reconstructions. The \emph{statistical
outlier filter} computes for every point the mean distance $\bar d_i$ to its $k$ nearest
neighbours, and removes the points with $\bar d_i>\mu+\alpha s$, where $\mu$ and $s$ are the mean
and standard deviation of all $\bar d_i$. The eigenvalues of the local covariance also carry
geometric information. The \emph{surface variation} \citep{Pauly2002}
\begin{equation}
\sigma_\lambda=\frac{\lambda_0}{\lambda_0+\lambda_1+\lambda_2}\in\bigl[0,\tfrac13\bigr]
\end{equation}
is zero on planes, grows with curvature and reaches $1/3$ for isotropic distributions; it behaves
like a curvature estimate and flags edges, corners and noisy regions (\cref{fig:outliers}).

\fig[0.85\textwidth]{ch09_outliers.png}{Left: the test surface with $6\%$ uniformly distributed
outliers (red). Centre: points flagged by the statistical filter ($k=16$, $\alpha=1$). Right:
surface variation $\sigma_\lambda$ on the clean surface, high on the saddle regions and the inner
side of the handles.}{fig:outliers}

\section{Why an implicit representation}

Reconstruction methods are either \emph{explicit}, connecting the input points directly into
triangles (ball pivoting \citep{Bernardini1999}, Delaunay-based methods, alpha shapes), or
\emph{implicit}, fitting a function $f:\R^3\to\R$ whose zero level set approximates the points and
extracting it as a mesh. Explicit methods interpolate the data, which is desirable for clean,
dense samples but propagates noise and leaves holes. Implicit methods approximate: they average
noise, produce watertight manifold surfaces by construction, fill holes, and their topology is
decided by the function rather than by local connectivity decisions (\cref{sec:implicit}). The rest
of the chapter follows the implicit route, which needs a way to extract the zero set.

\section{Extracting iso-surfaces}\label{sec:marching}

\subsection{Marching squares}

Consider first a 2D function $F$ sampled on a grid. Each cell has four corners, each inside
($F<0$) or outside ($F\ge0$), giving $2^4=16$ configurations. In each one the zero level crosses
the edges whose endpoints have different signs, and the crossing on edge $(\ba,\bb)$ is placed by
linear interpolation,
\begin{equation}
\bx=\ba+t(\bb-\ba),\qquad t=\frac{F(\ba)}{F(\ba)-F(\bb)} .
\end{equation}
A lookup table gives the segments for each configuration. Two configurations (diagonally
opposite corners inside) are ambiguous: the two segments can separate the inside corners or join
them (\cref{fig:msquares}). Evaluating the bilinear interpolant at the centre of the cell (the
\emph{asymptotic decider} \citep{Nielson1991}) chooses the consistent option. Linear
interpolation makes the extracted curve accurate even on coarse grids: a circle of radius one
extracted on a $9\times9$ grid has length $6.228$ instead of $2\pi=6.283$.

\fig[0.85\textwidth]{ch09_marching_squares.pdf}{Marching squares. From left: a slice $z=0$ of the
signed function of the test surface with its zero level; a unit circle extracted on a $9\times9$
grid (filled dots inside); the two possible connections of the ambiguous case 5.}{fig:msquares}

\subsection{Marching cubes and marching tetrahedra}

Marching cubes \citep{Lorensen1987} applies the same idea to cubes: $2^8=256$ sign
configurations, which reduce to 15 by symmetry, each mapped to up to four triangles with vertices
interpolated on the edges. The original table had ambiguous faces whose choices were not consistent
between neighbouring cubes, producing holes; corrected tables and the asymptotic decider fix this.
\emph{Marching tetrahedra} splits each cube into five or six tetrahedra. A tetrahedron has only
$2^4=16$ configurations of three kinds: no crossing, one corner separated from the others (one
triangle), or two corners separated from two (a quadrilateral, two triangles). There is no
ambiguity, at the price of more triangles and a result that depends on how the cubes are split.
Both methods produce meshes with many thin triangles, which remeshing improves (\cref{ch:remesh}).

\section{Hoppe's signed distance}

\Citet{Hoppe1992} define the signed distance to the surface through the tangent planes estimated
from the points. Each point $\bp_i$ has a tangent plane with centre $\bo_i$ (the centroid of its
neighbours) and oriented normal $\bn_i$. For a query $\bx$, find the plane whose centre is closest
and use the signed distance to it:
\begin{equation}
f(\bx)=(\bx-\bo_i)\cdot\bn_i,\qquad i=\argmin_j\norm{\bx-\bo_j}.
\end{equation}
Where $\bx$ is far from all the data the function is left undefined, so holes remain open. The
function is only piecewise smooth (it jumps across the Voronoi boundaries of the centres), and
noise in the normals produces spurious sign changes, small handles and isolated components.

\section{Radial basis functions}

\Citet{Carr2001} fit a smooth global interpolant
\begin{equation}
f(\bx)=\sum_{j=1}^Nw_j\,\varphi(\norm{\bx-\bc_j})+P(\bx),
\end{equation}
with the \emph{biharmonic} kernel $\varphi(r)=r$ (the minimiser of the bending energy
$\int\norm{\grad^2f}^2$ in 3D) and a linear polynomial $P$. The interpolation conditions require
$f=0$ at the input points. Since $f\equiv0$ also satisfies them, \emph{off-surface} points are added
along the normals: $f(\bp_i+\varepsilon\bn_i)=\varepsilon$ and $f(\bp_i-\varepsilon\bn_i)=-\varepsilon$.
With centres $\bc_j$ at all constraint points, the weights and polynomial coefficients solve the
symmetric system
\begin{equation}
\begin{pmatrix}\mA&\mP\\\mP\T&\bzero\end{pmatrix}\begin{pmatrix}\bw\\\bc\end{pmatrix}=
\begin{pmatrix}\bfun\\\bzero\end{pmatrix},\qquad A_{ij}=\varphi(\norm{\bc_i-\bc_j}),\quad
\mP_{i\cdot}=(1,\bc_i\T),
\end{equation}
whose second block row enforces the orthogonality conditions that make the solution unique. The
matrix is dense, so a direct solve costs $O(N^3)$; Carr et al.\ use fast multipole methods and a
greedy selection of the centres to handle millions of points. Replacing $\mA$ by $\mA+\rho\mI$
turns interpolation into smoothing, which is needed for noisy data. RBF surfaces are smooth and
fill holes gracefully, since the biharmonic function extends the surface with minimal bending.

\section{Poisson reconstruction}

\subsection{From normals to an indicator function}

\Citet{Kazhdan2006} reconstruct the \emph{indicator function} $\chi$ of the solid, equal to one
inside and zero outside. Its gradient is zero everywhere except at the surface, where it is a delta
function in the direction of the inward normal. After smoothing with a filter $\tilde F$, the
gradient of the smoothed indicator equals the smoothed field of (inward) normals:
\begin{equation}
\grad(\chi*\tilde F)(\bq)=\int_{\partial M}\tilde F_\bp(\bq)\,\bn_{\mathrm{in}}(\bp)\dd\bp\approx
\sum_{\bp\in P}A_\bp\,\tilde F_\bp(\bq)\,\bn_{\mathrm{in}}(\bp)=:\bm V(\bq),
\end{equation}
where the integral is approximated by the samples with area weights $A_\bp$. The oriented points
therefore give us the gradient of the function we want. In general $\bm V$ is not exactly a gradient
field, so we look for the $\chi$ whose gradient is closest in the least-squares sense,
\begin{equation}
\min_\chi\int\norm{\grad\chi-\bm V}^2\dd\bx\quad\Longrightarrow\quad\Lap\chi=\grad\cdot\bm V ,
\end{equation}
a Poisson equation, by the Euler--Lagrange equation of the functional.

\subsection{Solving and extracting}

On a regular grid with periodic boundary conditions the Poisson equation is diagonal in the
Fourier basis, so the notebook solves it with the FFT: transform $\grad\cdot\bm V$, divide each
coefficient by $-\norm{\boldsymbol\omega}^2$ (setting the constant mode to zero), and transform back.
The original method uses an adaptive octree with B-spline basis functions centred at the nodes and
a multigrid solver, so the resolution concentrates near the surface. The solution is defined up to
a constant, so the surface is extracted at the iso-value equal to the average of $\chi$ at the
input samples, with marching cubes. The method is global, so it is robust to noise and missing
data, and produces watertight surfaces; its weaknesses are over-smoothing of thin features and the
tendency to close holes with smooth caps whether or not they are real.

\subsection{Screened Poisson}

\Citet{Kazhdan2013} add a \emph{screening} term that asks the function to vanish at the samples
(taking the iso-value $\tfrac12$ into account, or equivalently shifting $\chi$),
\begin{equation}
\min_\chi\int\norm{\grad\chi-\bm V}^2\dd\bx+\alpha\sum_{\bp\in P}\chi(\bp)^2 ,
\end{equation}
whose optimality condition is the screened Poisson equation $(\Lap-\alpha\mathcal I_P)\chi=\grad\cdot\bm V$,
with $\mathcal I_P$ the operator that evaluates at the samples. The interpolation weight $\alpha$
pulls the surface towards the data and recovers detail lost by the plain method, while keeping
its robustness. The octree depth controls the resolution (\cref{fig:bunny-poisson}).

\fig{ch09_bunny_poisson.png}{Screened Poisson reconstruction (Open3D implementation) of six merged
range scans of the bunny at octree depths 5, 7 and 9. Higher depth resolves more detail and also
more noise.}{fig:bunny-poisson}

\section{Comparison}

\Cref{fig:recon} compares the three implicit methods on the noisy test surface, with error maps
against the exact surface. All recover the correct topology ($\chi=0$, genus one), and the median
errors are similar: $0.077\%$, $0.073\%$ and $0.091\%$ of the bounding-box diagonal for Hoppe, RBF
and Poisson (95th percentiles $0.234\%$, $0.230\%$ and $0.251\%$). The differences appear in harder
regimes (\cref{fig:regimes}): with strong noise Hoppe's piecewise construction produces a surface of
genus 35, full of spurious handles, while RBF and Poisson remain genus one; with a large hole Hoppe
leaves it open and creates spurious handles at its border, while the global methods fill it; with
sparse data all three succeed on this simple shape.

\fig{ch09_reconstruction.png}{Left: a slice of the indicator function computed by Poisson
reconstruction (the dashed line is the extracted iso-value). Right: error maps of Hoppe, RBF and
Poisson reconstructions on the test surface (dark is small error, bright is large).}{fig:recon}

\fig[0.85\textwidth]{ch09_regimes.png}{The three methods in three difficult regimes: strong noise
($\sigma=0.02$), a large hole, and only 1200 points. The genus of each result is shown in the
title; the correct value is one.}{fig:regimes}

\paragraph{Beyond implicit functions.} Man-made objects are better reconstructed by detecting
primitives first and structuring the points around them \citep{Lafarge2013}, or by fitting
primitives with global relations \citep{Li2011}. Learning-based methods add priors learned from
data: neural implicit fields encode $f$ as a network, and Point2Mesh \citep{Hanocka2020} shrink-wraps
an initial mesh onto the points while a convolutional network over the mesh acts as a
\emph{self-prior} that favours locally self-similar geometry.

\begin{notebook}{Tema07\_Normales\_y\_Reconstruccion.ipynb}
Block 1: total least squares plane fitting, vectorised PCA normals, the neighbourhood sweep and
$k^\star(\sigma)$, orientation by minimum spanning tree, the statistical outlier filter and surface
variation. Block 2: marching squares, the marching tetrahedra case table, Hoppe's signed distance,
biharmonic RBFs with off-surface points, Poisson reconstruction with an FFT solver and the
iso-value, and the comparison of the three methods in the three regimes of \cref{fig:regimes}.
\end{notebook}

\begin{keypoints}
\begin{itemize}[leftmargin=*]
  \item PCA normals are the eigenvectors of the smallest eigenvalue of the local covariance; the
  neighbourhood size trades noise against curvature.
  \item Orientation is propagated over a minimum spanning tree with weights $1-\abs{\bn_i\cdot\bn_j}$.
  \item Implicit reconstruction fits $f$ and extracts $\{f=0\}$ with marching squares, cubes or
  tetrahedra, interpolating linearly on the edges.
  \item Hoppe uses the tangent plane of the closest sample (local, fragile); RBFs solve a dense
  global interpolation with off-surface constraints; Poisson solves $\Lap\chi=\grad\cdot\bm V$ for the
  indicator function, and its screened version pulls the surface to the points.
\end{itemize}
\end{keypoints}

\section*{Exercises}
\begin{exercise}
Show that the centroid is the optimal plane centre in total least squares and that the minimum of
the objective equals the smallest eigenvalue of $\mC$.
\end{exercise}
\begin{exercise}
Write the 16 cases of marching tetrahedra grouped into the three types, and the triangles produced
in each, with consistent orientation.
\end{exercise}
\begin{exercise}
Derive the Euler--Lagrange equation of $\int\norm{\grad\chi-\bm V}^2$ and explain why the solution is
determined only up to a constant. How is the constant fixed in practice?
\end{exercise}
